\section{本章小结}
\label{sec:pdenum-summary}

开头的问题是：怎样把连续的场换成有限个数，这些数又代表什么。三种方法给出了三种回答，它们分别保留了连续问题的不同结构，表~\ref{tab:pdenum-methods}作了对比。

\begin{BookInlineTable}
 \centering
 \caption{三种离散方法的对比。}
 \label{tab:pdenum-methods}
 \begin{tabular}{p{16mm}p{20mm}p{20mm}p{38mm}}
  \toprule
  方法 & 未知量 & 保留的结构 & 精度保证依赖\\
  \midrule
  有限差分 & 节点值 & 局部导数 & 解的光滑性，离散逆算子的界\\
  Galerkin有限元 & 基函数系数 & 弱形式 & 有界、强制、协调，解的正则性，网格质量\\
  守恒有限体积 & 单元平均 & 局部守恒 & 守恒、相容，熵选择，全变差控制\\
  \bottomrule
 \end{tabular}
\end{BookInlineTable}

有限差分用Taylor展开逼近导数。定理~\ref{thm:pdenum-fd-error}说明，$u\in C^4$时小的截断误差经由与$h$无关的$\lVert\bm{A}_h^{-1}\rVert_\infty$转成小的解误差；解不够光滑时，阶数会下降。扩散问题上，网格加密带来更快的衰减模式，稳定性和精度对时间步提出了不同的要求。

Galerkin方法把弱问题限制到有限维子空间。Céa估计把误差化成最佳逼近问题，只用到有界性、强制性和协调性。非齐次Dirichlet数据要作提升，右端的边界贡献不能丢；实际的收敛阶还要结合解的正则性、单元次数和网格质量。

有限体积方法用界面通量更新单元平均，内部通量抵消保证了离散守恒，也保证了激波速度的正确。但守恒只管总量：Roe型通量相容又守恒，却会保留违反熵条件的膨胀激波，需要单调性或熵不等式来选出正确的解。

收敛的论证在线性和非线性问题上走了不同的路。线性格式由一致性和稳定性得到收敛，Lax等价定理说明两者在适定问题上等价；Fourier分析只适用于规则网格、常系数和周期边界，边界闭合要另行检查；CFL条件是依赖区域的必要条件，并不充分。非线性守恒格式则分别核验紧性、Lax--Wendroff定理的弱解极限和离散熵不等式，不能套用线性的结论。

因此，得到一个离散解只是开始。要说明它代表了什么连续解，还需要把空间、时间、积分和代数误差分开检验，用网格加密观察实际的收敛阶，并确认解的尺度已被网格分辨。
